%%%%%%%%%%%%%%%%%%%%%%%%%%  example_doublecolumn.tex  %%%%%%%%%%%%%%%%%%%%%%%%%%%%%%%%
\documentclass[twoside,leqno,twocolumn]{article}

% Comment out the line below if using A4 paper size
\usepackage[letterpaper]{geometry}

\usepackage{siamproceedings}

\usepackage[T1]{fontenc}
\usepackage{amsfonts}
\usepackage{graphicx}
\usepackage{epstopdf}
\usepackage{enumitem}
\usepackage{algorithmic}
\ifpdf
  \DeclareGraphicsExtensions{.eps,.pdf,.png,.jpg}
\else
  \DeclareGraphicsExtensions{.eps}
\fi

\newcommand{\creflastconjunction}{, and~}
\newsiamremark{remark}{Remark}
\newsiamremark{hypothesis}{Hypothesis}
\crefname{hypothesis}{Hypothesis}{Hypotheses}
\newsiamthm{claim}{Claim}

\usepackage{amsopn}
\DeclareMathOperator{\diag}{diag}

\usepackage{amsmath,amssymb}
\usepackage{amsthm}
\usepackage{algorithm}
\usepackage{algpseudocode}
\usepackage{booktabs}
\usepackage{graphicx}
\usepackage{subcaption}
\usepackage[hidelinks]{hyperref}
\usepackage{color}

\graphicspath{
  {./}
  {./results/}
  {./results/multidata/}
  {./OverclusterDRG_Paper/}
  {./OverclusteringRKC/}
  {./overclustering_me/}
}

\newcommand{\OPT}{\mathrm{OPT}}
\newcommand{\rz}{r_z}
\newcommand{\CL}{C_L}
\newcommand{\Copt}{C^{*}_{Q}}
\newcommand{\rkc}{R$k$C}

\begin{document}

\newcommand\relatedversion{}

\title{\Large OverclusterDRG: a Deterministic, LP-Free Heuristic for Robust $k$-Center\relatedversion}

\date{}

\maketitle

\fancyfoot[R]{\scriptsize{Copyright \textcopyright\ 20XX by SIAM\\
Unauthorized reproduction of this article is prohibited}}

\begin{abstract}
...
\end{abstract}


\section{Introduction}

The $k$-center problem is a foundational framework in metric clustering and facility location. Given a dataset and a budget of $k$ centers, the objective is to place the centers such that the maximum distance from any data point to its nearest center is minimized. Because it optimizes a worst-case distance, $k$-center is the standard formulation for applications requiring strict service guarantees, such as emergency facility placement, sensor coverage, and bounding network latency. While the exact problem is NP-hard, Gonzalez \cite{gonzalez1985} introduced an elegant $O(Nk)$-time greedy algorithm that achieves a tight 2-factor approximation. The algorithm operates via a farthest-point traversal: starting with an arbitrary point, it iteratively selects the point farthest from the current set of centers.

However, in the rigid min-max objective of the classical $k$-center problem, a single distant outlier, measurement error, or anomalous data entry can dictate the objective. It can drag center away from dense, meaningful clusters merely to cover the isolated point. To handle this, the Robust $k$-Center problem (\rkc) introduces an outlier budget $z \ge 0$. The algorithm is allowed to discard up to $z$ points from the dataset, and the objective is evaluated only on the remaining $N - z$ non-outlier points. While this robust formulation protects the clustering from noise, it introduces a significant computational burden: identifying the optimal subset of outliers superimposes a combinatorial selection problem on top of an already NP-hard clustering task.


While \rkc~elegantly models real-world noise, algorithms for solving it reveal a severe disconnect between theoretical worst-case guarantees and practical algorithm engineering. A natural practitioner baseline is to run the Gonzalez traversal for $k$ steps and evaluate the radius after discarding the $z$ farthest points. Unfortunately, all theoretical guarantees vanish under this approach: a single outlier can irreparably hijack a center selection, leaving true clusters entirely unrepresented. Consequently, the theoretical literature has developed approximation algorithms with strict formal guarantees, but these algorithms present computational bottlenecks at scale:
\begin{itemize}
    \item \textbf{Combinatorial Approximation:} Charikar et al. \cite{charikar2001} introduced the first constant-factor approximation for robust $k$-center, yielding a 3-factor approximation. While theoretically significant, their algorithm relies on optimal radius-guessing, constructing a graph of intersecting disks, and greedily selecting expanded clusters. Our empirical profiling reveals that this method scales as $O(N^{2.12})$ in practice, requiring hundreds of seconds on moderately large datasets.
    \item \textbf{LP rounding schemes:} The approximation factor for \rkc~was later improved to a tight 2 using linear programming (LP) rounding by Chakrabarty, Goyal, and Krishnaswamy \cite{}; Harris et al. \cite{}. Exact formulations using Linear Programming (LP) are even worse; our evaluations show that solving an exact LP requires up to 94 hours for a single configuration at $N=20,000$.
    \item \textbf{Randomized sampling:} To bypass these computational bottlenecks, Ding, Yu, and Wang \cite{ding2019} proposed a randomized greedy strategy that achieves a 2-approximation with a constant probability per trial. To guarantee a 99\% probability of success, their algorithm requires running $T = \lceil(z+1) \ln 100\rceil$ independent randomized trials.
    Because the outlier budget $z$ is often a constant fraction of $N$ in real-world applications this trial multiplier dominates the runtime. For example, on a dataset of 20,000 points with a 20\% outlier rate, Ding's algorithm requires 18,423 redundant passes, taking 431 seconds to complete. The runtime scales linearly with the outlier budget, forcing practitioners to abandon both determinism and speed.
   \item  \textbf{Naive Heuristics:} Given that algorithms with provable guarantees are not practical, a natural fast heuristic is a simple adaptation of Gonzalez's algorithm: \textcolor{red}{run the standard farthest-point traversal for $k$ steps to select candidate centers, and then assign points to their nearest center while discarding the $z$ points that contribute most to the cost.} Post-hoc removal of the $z$ farthest points operates in fast $O(N(k+z))$ time but offers no formal guarantees and yields poor quality on real-world datasets; degrading to 1.28x to 1.92x the optimal LP radius in our evaluations.
\end{itemize}

There remains a need for an algorithm that provides the empirical quality of rigorous combinatorial and LP-based methods, operates deterministically, and scales to large datasets without prohibitive computational overhead.


In this paper, we present \textbf{OverclusterDRG} - \textbf{Over}clustering
followed by a \textbf{D}eterministic \textbf{R}obust \textbf{G}reedy
selection; a deterministic, LP-free heuristic for \rkc. We leverage the theoretical insight of candidate oversampling - a standard data reduction technique in the parameterized complexity (FPT) literature - but replace its intractable subset-guessing step with a highly optimized, deterministic heuristic pipeline. By combining farthest-point over-sampling with a local search mechanism, we provide a practical method that scales well empirically while maintaining consistent solution quality. Our algorithm operates in two phases: it first deterministically builds a candidate pool of $k+O$ points using a permutation-invariant Gonzalez farthest-point rule (where $O$ is an overcount parameter). Outliers will still attract the farthest-point rule, but with $O$ spare steps, the pool can absorb them and still retain a representative center for every true cluster. The right amount of slack turns out to be exactly $O=z$. It then extracts $k$ centers using a forward-greedy initialization (OC-Forward) followed by a 1-swap local search (OC-Local).

Our primary contributions are:
\begin{enumerate}
    \item A Deterministic, Single-Pass Framework: we bypass the randomized trial multipliers present in recent scalable algorithms. By utilizing strict lexicographical tie-breaking, our Phase 1 pool construction is entirely deterministic and permutation-invariant. OverclusterDRG requires only a single pass to build its candidate pool, completing the hardest evaluated configuration ($N=20,000$, $z=20\%$) in just 57 seconds, avoiding the 18,423 independent trials required by randomized equivalents.

    \item Empirical Quality Evaluated Against Exact LP Bounds: to baseline heuristic performance, we computed exact LP lower bounds across 102 configurations on three real-world datasets. OverclusterDRG achieves a mean LP approximation ratio of 1.147 on the UCI Adult dataset, outperforming the achieved radii of both Charikar's (1.173) and Ding's (1.235) implementations. This ordering replicates across secondary datasets (1.155 on Diabetes, 1.171 on Covertype). Furthermore, in exhaustive evaluations on 400 small subsamples, our 1-swap local search perfectly recovered the exact combinatorial pool optimum in 97.2\% of instances, with a maximum observed gap of only 4.82\%.
\item  Computational Scalability: on the Adult dataset, OverclusterDRG completes with a median runtime of 0.36 seconds, yielding a median 7x speedup over Charikar (3.60s) and an 8x speedup over Ding (2.22s). On a configuration-by-configuration basis, it attains a tighter radius than Charikar in 33 of 36 cases, and outperforms Ding in 35 of 36 cases. Utilizing memory-efficient distance matrix caching, its empirical time complexity scales at roughly $O(N^{1.85})$. At large scales ($N = 99,000$), where exact LPs and Charikar's algorithm are computationally intractable, OverclusterDRG remains efficient, completing a pass in 1,100 seconds---more than three times faster than Ding's randomized baseline (3,800 seconds).
\end{enumerate}

\section{Related Work}

Gonzalez \cite{gonzalez1985} gave the farthest-point heuristic, a
$2$-approximation for $k$-center without outliers in $O(Nk)$ time. Its
\rkc~adaptation evaluates the resulting centers under $\rz$ and is the
fastest baseline we test, but the guarantee evaporates: a single
outlier hijacks a center, and no bound on $\rz$ survives.

Charikar et al.\ \cite{charikar2001} introduced the first
constant-factor approximation for clustering with outliers; the greedy
disk-based variant for RKC certifies $3r^{*}$, where $r^{*}$ is the
smallest feasibility radius found. Quality is good. The cost is a
quadratic-in-$N$ candidate-radius search that becomes painful past a
few thousand points.

Ding, Yu, and Wang \cite{ding2019} showed that a randomised greedy
strategy succeeds with constant probability per trial, amplified to
$99\%$ with $T=\lceil(z{+}1)\ln 100\rceil$ repetitions. The method
scales gently in $N$ but the multiplier $T$ grows with $z$: at
$z=20\%$ of $N=20{,}000$ this is $18{,}423$ trials. The output is also
random---two runs disagree.

Local search underpins the best combinatorial bounds for $k$-median:
Arya et al.\ \cite{arya2004} proved a $(3+2/p)$-approximation for
$p$-swap local search. Their telescoping argument sums local
optimality conditions over swaps, which works because $k$-median is a
sum of distances. The $k$-center objective is a maximum, and the
argument breaks; this failure is precisely what makes our locality-gap
conjecture (Section~\ref{sec:conjecture}) open. On the parameterised
side, Goyal and Jaiswal \cite{goyal2021} give tight FPT approximations
for constrained $k$-center variants; these are exponential in $k$ and
target a different regime.

No prior algorithm is simultaneously deterministic, LP-free, fast at
large $N$ and large $z$, and formally tied to $\OPT$. That is the gap
this paper fills.


\section{Problem Setup}
\label{sec:setup}

Let $(X,d)$ be a finite metric space with $|X|=N$. Given integers
$k\ge 1$ and $z\ge 0$, find $C\subseteq X$ with $|C|=k$ and
$S\subseteq X$ with $|S|\le z$ minimizing
\[
  R(C,S) \;=\; \max_{x\in X\setminus S} d(x,C),
  \qquad d(x,C)=\min_{c\in C} d(x,c).
\]
The optimum is $\OPT = \min_{|C|=k,\,|S|\le z} R(C,S)$.

In simple words, place $k$ centers and delete up to $z$ points so that every surviving point sits close to a center; minimize the worst surviving distance.

\paragraph{(Linear Program) LP.}
For a fixed radius $r$, define $N(j,r)=\{i \mid d(i,j)\le r\}$ as the
set of points close enough to serve as a center for $j$.
Introduce variables: $y_i\in[0,1]$ (center opened at $i$);
$x_{ij}\in[0,1]$ (point $j$ assigned to center $i$);
$c_j\in[0,1]$ (point $j$ is covered, i.e.\ not an outlier).
\begin{align*}
\sum_{i\in N(j,r)} x_{ij} &\ge c_j \quad\forall j
  &&\text{($j$ covered only if assigned within radius $r$)}\\
x_{ij} &\le y_i \quad\forall i,j
  &&\text{(assign only to open centers)}\\
\sum_i y_i &\le k
  &&\text{(at most $k$ centers)}\\
\sum_j c_j &\ge N-z
  &&\text{(at most $z$ outliers)}\\
x_{ij},\,y_i,\,c_j &\in [0,1]
\end{align*}
We binary-search over all pairwise distances as candidates for $r$ and
solve this feasibility problem at each probe. The smallest feasible $r$
is $R_{\mathrm{LP}}$, which satisfies $R_{\mathrm{LP}}\le\OPT$ and
serves as the lower bound for all reported ratios.
The LP is expensive - beyond $N\approx 2{,}000$ a single configuration costs
hours to days - and we pay that cost in full: every configuration of
every dataset uses an exact LP, solved to a feasibility tolerance of
$10^{-6}$, up to 94 hours at $N=20{,}000$.

\textcolor{red}{The analysis compares the \rkc~optimum against a standard $k$-center
optimum with no outliers:
\[
  \OPT_{k,z} = \OPT, \qquad
  \OPT_{m,0} = \min_{|C|=m}\;\max_{x\in X} d(x,C).
\]
The first discards $z$ points; the second covers all $N$ points with
$m$ centers. Their relationship for $m=k+O$ is the hinge of the entire
theory.}


\section{The OverclusterDRG Algorithm}
\label{sec:algorithm}

The algorithm has one idea and three phases. The idea: Gonzalez wastes
steps on outliers, so give it extra steps. The phases: build an
overcomplete pool, initialise greedily, polish with $1$-swap local
search. The name records the ingredients---\textbf{Over}clustering
followed by a \textbf{D}eterministic \textbf{R}obust \textbf{G}reedy
selection.

\paragraph{Phase 1: pool construction.}
Run the Gonzalez farthest-point rule for $k+O$ steps. The first point
is the data point nearest the empirical centroid $\bar{X}$, which
makes the whole algorithm deterministic; every later point is the
point farthest from the current pool,
$q_{i+1}=\arg\max_x \min_{j\le i} d(x,q_j)$. Ties in both argmins
and argmaxes are broken lexicographically by point coordinates, so the
output is independent of input order and the determinism claim is
literal. We verified this empirically: re-running every configuration
under random input permutations produced byte-identical pools, so the
lexicographic rule is permutation-invariant in floating-point as well
as in exact arithmetic (the coordinate comparison, not the
floating-point distance, resolves ties). The pool is
$Q=\{q_1,\dots,q_{k+O}\}$ and costs $O(N(k{+}O))$ to build. The
default overcount is $O=z$, the outlier budget---the choice
Section~\ref{sec:theory} proves correct.

\paragraph{Phase 2a: forward greedy (OC-Forward).}
Pre-compute the $N\times(k{+}O)$ distance matrix $D[x][i]=d(x,q_i)$
once. Starting from $\{q_1\}$, repeatedly add the pool point whose
inclusion minimises $\rz(C\cup\{q\},X)$; ties are broken by lowest
pool index. Each candidate evaluation is $O(N)$ using $D$ and a
running minimum, so the phase costs $O(NkO)$. The output $C_0$ is
already a strong solution and we report it separately as OC-Forward.

\paragraph{Phase 2b: $1$-swap local search (OC-Local).}
For each pair $(c_i,q)\in C\times(Q\setminus C)$, scanned in fixed
pool-index order, accept the first swap that strictly decreases
$\rz$, and repeat until no swap improves. Each sweep costs $O(NkO)$;
the objective strictly decreases and there are finitely many
$k$-subsets of $Q$, so termination is guaranteed. In practice
$5$--$10$ accepted swaps suffice.

\begin{algorithm}[t]
\caption{Gonzalez$(X, k)$}
\label{alg:gonzalez}
\begin{algorithmic}[1]
\State Start from an arbitrary point. Set $C = \{c_1\}$.
\State Add the point farthest from all current centers in $C$.
\State Repeat Step 2 until $k$ centers are selected.
\State \Return $C$
\end{algorithmic}
\end{algorithm}

\begin{algorithm}[t]
\caption{OverclusterDRG$(X, k, z, O{=}z)$}
\label{alg:main}
\begin{algorithmic}[1]
\State $Q \gets \textsc{Gonzalez}(X,\, k+O)$
  \Comment{Phase 1: pool of $k{+}O$ candidates}
\State $C_0 \gets \textsc{OC-Forward}(X, Q, k, z)$
  \Comment{Phase 2a: greedy initialisation}
\State $C \gets \textsc{OC-Local}(X, Q, C_0, z)$
  \Comment{Phase 2b: $1$-swap improvement}
\State \Return $C$
\end{algorithmic}
\end{algorithm}

\paragraph{Complexity.}
Phase~1 costs $O(N(k{+}O))$. Each Phase-2 sweep---one full pass over
the $k(k{+}O)$ candidate swaps, each evaluation $O(N)$ via the
precomputed matrix---costs $O(NkO)$, which is $O(Nkz)$ at the default
$O=z$. The number of sweeps is bounded by the number of distinct
$k$-subsets of $Q$ and is therefore finite, but we know no polynomial
worst-case bound; empirically $5$--$10$ accepted swaps and a comparable
number of sweeps suffice on every configuration we ran. We therefore
state the cost as $O(Nkz)$ \emph{per sweep} and report the observed
sweep count, rather than claim a worst-case total. When $z$ is a
constant fraction of $N$, as in our entire grid, the per-sweep cost is
$O(N^{2}k)$---the same asymptotic order as Charikar's measured
$N^{2.12}$, and consistent with our own measured exponent
$O(N^{1.85})$ ($R^2=0.976$). The advantage over Charikar is a large
constant factor, not an order-of-growth separation; the advantage over
Ding is structural---one deterministic pass in place of
$\lceil(z{+}1)\ln 100\rceil$ randomised trials.

\emph{Memory.}
The $N\times(k{+}O)$ distance matrix dominates: $O(Nz)$ doubles at
$O=z$, about $7.9$\,GB at $N=99{,}000$, $k=20$, $z=10\%$, which is the
largest configuration that fits in core on our machine and sets the
practical single-machine ceiling for full matrix residence. At
$k=50$, $z=20\%$, $N=10^{5}$ the matrix would need $\approx8$\,TB, far
beyond RAM; Phase~2 then processes the matrix in row blocks, holding
$B$ rows at a time and recomputing each block's distances per sweep.
This caps memory at $O(B(k{+}O))$ and multiplies Phase-2 time by at
most the (small, observed $5$--$10$) sweep count, since each row is
rebuilt once per sweep instead of cached---a constant-factor cost tied
to the sweep count, not to $N$. For the theoretical guarantee one
replaces Phase~2 by exhaustive search over all $\binom{k+O}{k}$ pool
subsets at $O(N)$ each, giving $O(Nk^{O})$---polynomial for constant
$O$, but astronomically large at the default $O=z$; the exhaustive
variant is a vehicle for the guarantee, not a practical mode
(Theorem~\ref{thm:exhaustive}).


\section{Experiments}
\label{sec:experiments}

\subsection{Setup}
\label{sec:expsetup}

\begin{table}[t]
\caption{Datasets. All features z-score standardised.}
\label{tab:datasets}
\centering
\begin{tabular}{lrrl}
\toprule
Dataset & Rows used & Features & Source \\
\midrule
UCI Adult & 32{,}561 & 6 & \cite{adult_dataset} \\
Diabetes Indicators & 99{,}492 & 8 & UCI \\
Covertype & 30{,}000 of 581{,}012 & 10 & \cite{covertype_dataset} \\
HTRU2 & 17{,}898 & 8 & \cite{htru2} \\
\bottomrule
\end{tabular}
\end{table}

Table~\ref{tab:datasets} lists the datasets. The Adult grid is
$N\in\{500,1000,2000,5000,10000,20000\}$, $k\in\{10,20,50\}$, and
outlier rate $z/N\in\{10\%,20\%\}$---36 configurations. Diabetes
substitutes $N=15{,}000$ for $2{,}000$ (36 configurations); Covertype
uses $N\in\{500,1000,5000,10000,15000\}$ (30 configurations). Five
algorithms compete: the RKC-adapted Gonzalez heuristic ($O(Nk)$, no
guarantee); OC-Forward (Phase~2a only); OC-Local (the full
algorithm, $O=z$); Charikar et al.\ \cite{charikar2001}, reported at
its achieved solution radius; and Ding et al.\ \cite{ding2019} at
$T=\lceil(z{+}1)\ln 100\rceil$ trials for a $99\%$ success guarantee.
The evaluation metric is the ratio of the achieved radius to the exact
LP lower bound of Section~\ref{sec:setup}; lower is better and
$1.000$ is unattainable in general.

\begin{algorithm}[t]
\caption{RKC-Gonzalez$(X, k, z)$}
\label{alg:rkcgonzalez}
\begin{algorithmic}[1]
\State Start from an arbitrary point. Set $C = \{c_1\}$.
\State Add the $(z{+}1)$-th farthest point from current centers,
       skipping the $z$ most extreme candidates.
\State Repeat Step 2 until $k$ centers are selected.
\State Drop the $z$ points farthest from their assigned center.
\State \Return $C$ and its radius.
\end{algorithmic}
\end{algorithm}

\paragraph{Preprocessing.}
Adult uses its six numeric attributes (age, fnlwgt, education-num,
capital-gain, capital-loss, hours-per-week) under Euclidean distance
after z-scoring; categorical attributes are dropped. Covertype is
subsampled to $30{,}000$ of $581{,}012$ rows by uniform random
sampling without replacement (NumPy \texttt{default\_rng}, seed $42$),
then z-scored; this single fixed sample is reused across all Covertype
configurations. At full size, Charikar's measured $N^{2.12}$ scaling
implies hundreds of hours per configuration and Ding's certificate
trial count is out of reach, which is why the subsample is used.

\paragraph{Exact LP bounds.}
The LP lower bound is solved exactly at every configuration of every
dataset (Gurobi 12.0.3, feasibility tolerance $10^{-6}$). Because the
optimal RKC radius is always a pairwise distance, the returned
$R_{\mathrm{LP}}$ is an exact pairwise distance, not an interpolant
(Appendix~\ref{app:lp}). Representative runtimes range from
$100$--$1{,}096$\,s at $N=2{,}000$ to $3.4\times10^{5}$\,s (94 hours)
at $N=20{,}000$. Because every algorithm within a configuration shares
the same exact denominator, all paired comparisons and orderings are
exact.

\paragraph{Free certificates in practice.}
Lemma~\ref{lem:cert} yields the certified bound $\OPT\ge
D_{k+z+1}/2$ from one extra Gonzalez step at every configuration. The
honest empirical verdict: the certificate is valid but conservative.
Across the three datasets the certified ratio
$2\,\rz(C)/D_{k+z+1}$ for OC-Local averages $3.59$ (Adult), $3.13$
(Diabetes), and $3.25$ (Covertype), and the certified lower bound
captures $36$--$39\%$ of the LP bound on average (range $0.16$--$0.62$).
The looseness has a structural cause: with $z=\Theta(N)$, the spacing
$D_{k+z+1}$ is the $(k{+}z{+}1)$-th farthest-point gap, which shrinks
well below $\OPT$ on dense data. The certificate is tightest at small
$N$ and small $z$ (ratio ${\approx}2$ at $N=500$) and is best read as a
free sanity bound, not an LP replacement; tight ratio reporting still
requires the LP. Appendix~\ref{app:lp} states the LP in full.

\subsection{UCI Adult: full comparison}

\begin{table}[t]
\caption{Aggregate LP ratios and runtimes over 36 UCI Adult
configurations. Ratio = radius / exact LP radius; lower is better.}
\label{tab:adult}
\centering
\begin{tabular}{lcccccc}
\toprule
Algorithm & Mean & Std & Min & Max & Med.\ time (s) & Max time (s) \\
\midrule
LP oracle & 1.000 & --- & 1.000 & 1.000 & 4{,}586 & 338{,}755 \\
\textbf{OC-Local} & \textbf{1.147} & 0.047 & 1.088 & 1.265 & 0.36 & 57 \\
OC-Forward & 1.158 & 0.047 & 1.104 & 1.283 & 0.26 & 52 \\
Charikar (achieved) & 1.173 & 0.056 & 1.104 & 1.306 & 3.60 & 564 \\
Ding & 1.235 & 0.032 & 1.179 & 1.331 & 2.22 & 431 \\
RKC-Gonzalez & 1.280 & 0.049 & 1.197 & 1.370 & 0.002 & 0.02 \\
\bottomrule
\end{tabular}
\end{table}

Table~\ref{tab:adult} summarises the headline comparison. OC-Local
attains the best mean ratio of every practical algorithm, $1.147$,
against Charikar's $1.173$, Ding's $1.235$, and Gonzalez's $1.280$. It
is also faster where it matters: median runtime $0.36$s versus
Charikar's $3.60$s and Ding's $2.22$s. Per configuration, OC-Local
beats Charikar in $33/36$ cases at a median $7\times$ speedup and beats
Ding in $35/36$ cases at a median $8\times$ speedup. The LP row is an
oracle, not an algorithm: its median runtime is $4{,}586$ seconds and
its worst case exceeds 94 hours.

\begin{figure}[t]
\centering
\includegraphics[width=0.95\linewidth]{fig3_quality_speed.png}
\caption{Quality--runtime tradeoff on UCI Adult; each point is one
$(N,k,z/N)$ configuration, runtime on a log axis, lower-left is
better. OC-Local (green diamonds) occupies the lower-left frontier at
both outlier rates: no other practical algorithm matches its radius at
any comparable runtime. LP (red stars) sits on the ratio-$1$ line at
runtimes orders of magnitude larger.}
\label{fig:qualityspeed}
\end{figure}

\begin{figure}[t]
\centering
\includegraphics[width=0.95\linewidth]{fig1_ratio_vs_N.png}
\caption{LP ratio versus dataset size $N$, averaged over $k$, one
panel per outlier rate. OC-Local (green) tracks the LP bound most
closely at every scale, with the gap to Gonzalez (grey, dashed) and
Ding (purple, dotted) widening rather than closing as $N$ grows.}
\label{fig:ratiovsN}
\end{figure}

Figure~\ref{fig:qualityspeed} shows the full tradeoff; OC-Local
dominates the lower-left corner. Figure~\ref{fig:ratiovsN} shows that
the advantage is stable in $N$: OC-Local's ratio stays near
$1.13$--$1.16$ from $N=500$ to $N=20{,}000$ at both outlier rates.

\begin{table}[t]
\caption{Paired tests on per-configuration ratio differences
($n=36$, exact LP throughout). Positive $\bar d$ means OC-Local is
better; the Wilcoxon column is the distribution-free signed-rank
test.}
\label{tab:ttests}
\centering
\begin{tabular}{lcccc}
\toprule
Comparison & $\bar d$ & $t$ & $p$ & Wilcoxon $p$ \\
\midrule
OC-Local vs RKC-Gonzalez & $+0.134$ & $+13.3$ & $<10^{-14}$ & $2.9\times10^{-11}$ \\
OC-Local vs Ding & $+0.089$ & $+14.3$ & $<10^{-15}$ & $5.8\times10^{-11}$ \\
OC-Local vs Charikar & $+0.026$ & $+5.7$ & $<10^{-5}$ & $3.1\times10^{-8}$ \\
OC-Local vs OC-Forward & $+0.012$ & $+7.9$ & $<10^{-8}$ & $2.9\times10^{-11}$ \\
\bottomrule
\end{tabular}
\end{table}

Table~\ref{tab:ttests} confirms that none of these gaps is noise. The
smallest comparison---OC-Local against its own initialisation
OC-Forward---is still significant at $p<10^{-8}$: the $1$-swap phase
improves the mean ratio by $0.012$ at a median cost of $0.10$ seconds,
and does so in 36 of 36 configurations. The polish is nearly free and
should always be applied. SimplifiedDRG, the pool-free deterministic
baseline, reaches only $1.295$; the pool buys an $11.5\%$ quality
improvement at $p<10^{-16}$.

The 36 configurations are slices of shared base data, so the paired
$t$-tests overstate independence; the Wilcoxon column gives the
distribution-free counterpart and agrees everywhere. As a further
robustness check, restricting to the 18 smallest configurations
($N\le 2{,}000$) leaves every conclusion intact: OC-Local wins
$17$--$18$ of $18$ against each baseline with Wilcoxon
$p\le 1.3\times10^{-4}$.

\paragraph{Runtime scaling.}
Log-log regression of runtime against $N$ gives OC-Local
$O(N^{1.85})$ ($R^2=0.976$), Charikar $O(N^{2.12})$ ($R^2=0.999$), and
Gonzalez $O(N^{0.80})$ ($R^2=0.939$); the sub-linear Gonzalez exponent
reflects fixed-overhead dominance at small $N$ and converges toward
$1$ at large $N$, consistent with $O(Nk)$ theory. Fitting the exponent
separately within each $(k,z/N)$ stratum, rather than pooling, leaves
OC-Local's $N$-exponent tightly clustered around $1.85$ (none below
$1.8$ or above $1.9$ across the six strata, all $R^2\ge0.95$), so the
pooled estimate is not an artifact of a particular $k$ or outlier rate.
Ding's nominal exponent is a gentle $O(N^{1.10})$, but the trial
multiplier $(z{+}1)\ln 100$ dominates in practice: at the hardest
configuration ($N=20{,}000$, $k=50$, $z=20\%$) Ding needs $18{,}423$
trials and $431$ seconds, while OC-Local completes its single
deterministic pass in $57$ seconds---a $7.5\times$ gap that the
multi-dataset experiments of Section~\ref{sec:multidata} widen further.

\subsection{Local search validation}
\label{sec:localvalid}

Theorem~\ref{thm:neg} says $1$-swap can fail arbitrarily badly. Does
it, on actual data? We compared OC-Local against exhaustive Phase~2 on
400 small instances. The instances are independent uniform random
subsamples of UCI Adult, each of size $N\le1000$ drawn without
replacement under a distinct seed, evaluated at $k\in\{5,10\}$ and
$z\le 10$; none is synthetic. This is the largest regime where
exhaustive search over $\binom{k+O}{k}$ subsets is feasible.

\begin{figure}[t]
\centering
\includegraphics[width=0.95\linewidth]{fig5_local_vs_exhaustive.png}
\caption{Gap between $1$-swap local search and exhaustive pool search
across 400 instances, plotted against pool size and coloured by $k$;
right panel: histogram of the 11 non-exact gaps. Local search is exact
in $97.2\%$ of instances and never exceeds the exhaustive optimum by
more than $4.82\%$.}
\label{fig:localexhaustive}
\end{figure}

OC-Local matched the exhaustive optimum exactly in 389 of 400
instances ($97.2\%$). The maximum observed gap was $4.82\%$, at the
largest pool (size 20, $k=10$, $z=10$), where exhaustive search
evaluates $184{,}756$ subsets in about 10 seconds and OC-Local
finishes in under 3 milliseconds. Every failure traces to the
adversarial structure of Theorem~\ref{thm:neg}---small clusters at
near-uniform distances---which random real-data slices essentially
never produce. Figure~\ref{fig:localexhaustive} shows the full gap
distribution. These 400 instances are also the empirical support for
the locality-gap conjecture: no observed gap exceeds $1+2/k$.

\subsection{Pool size analysis: the $O=z$ finding}
\label{sec:oablation}

Corollary~\ref{cor:phase} predicts a sharp behavioural change at
$O=z$. We sweep $O$ over 72 configurations on Diabetes and Covertype
and measure the radius of the best pool solution as a function of $O$.

\begin{figure}[t]
\centering
\includegraphics[width=0.95\linewidth]{fig_O_ablation.png}
\caption{Radius (normalised to its value at $O=z$) as a function of
pool overcount $O$ on a log axis, across all 72 Diabetes and Covertype
configurations. The curves fall steeply while $O<z$---the degradation
of Theorem~\ref{thm:necessity} made visible---and flatten near the
threshold, with the exact plateau onset varying by configuration.}
\label{fig:oablation}
\end{figure}

\begin{figure}[t]
\centering
\begin{subfigure}{0.48\linewidth}
  \includegraphics[width=\linewidth]{fig_phase_transition.png}
  \caption{}
  \label{fig:phasetransition}
\end{subfigure}
\hfill
\begin{subfigure}{0.48\linewidth}
  \includegraphics[width=\linewidth]{fig_O_vs_zk.png}
  \caption{}
  \label{fig:ovszk}
\end{subfigure}
\caption{(a) Radius versus $O$ for three representative configurations
with the $O=z$ threshold marked; quality degrades sharply left of the
line and is flat right of it. (b) The empirical knee
$O_{\mathrm{knee}}/z$ against $z/k$ across 72 configurations (Pearson
$r=-0.60$): more outliers per center push the knee toward $z$ itself,
while dimensionality shows no relationship.}
\label{fig:phasepair}
\end{figure}

Figure~\ref{fig:oablation} shows the sweep: radius drops steeply while
$O<z$ in every configuration. Define the knee $O_{\mathrm{knee}}$ as
the smallest $O$ within $2\%$ of the best achievable pool quality.
Across the 72 configurations the knee sits at $O_{\mathrm{knee}}/z$
between $1.0$ and $1.7$ with mean $1.55$; past $O=2z$, further gains
fall below $1\%$.

The theorem and the knee measure different quantities, and the
difference deserves stating plainly. Corollary~\ref{cor:phase}
concerns what the pool \emph{contains}, and containment transitions
exactly at $O=z$. The knee measures what OC-Local \emph{finds}, which
keeps improving modestly past $z$ before flattening at an average of
$1.55z$. The two agree on the steep left-of-$z$ degradation; they
differ on where improvement stops, and that gap between containment
and recovery is quantified by Table~\ref{tab:osettings} and is itself
an open question (Section~\ref{sec:conjecture}).
Theorem~\ref{thm:anytime} supplies the missing theory for the
post-threshold regime: the containment guarantee is
$(1+2\rho(O))\cdot\OPT$ and improves for as long as the Gonzalez
distances keep falling. Measured across the ablation configurations,
mean $\rho$ falls from $1.0$ at $O=z$ to $0.85$ at $O=1.5z$, $0.76$ at
$O=2z$, $0.63$ at $O=3z$, and $0.50$ at $O=5z$, tightening the
worst-case containment factor from $3\cdot\OPT$ to $2\cdot\OPT$ at
$O\approx5z$. The finer prediction---that the empirical knee coincides
with the flattening of $\rho$---is not settled by our data: under the
$2\%$ flattening rule only 4 of 72 configurations show a detectable
$\rho$-plateau within the swept range (correlation $0.69$ at $n=4$,
too few to conclude). The anytime bound holds; whether it explains the
knee is open.

\paragraph{Adaptive rule in practice.}
The stopping rule of Corollary~\ref{cor:adaptive} was run at
$\varepsilon\in\{1.0,\,0.5,\,0.25,\,0.1\}$. The $\varepsilon=1$ run is
an implementation sanity check rather than a result: the rule meets
its threshold at $t=z$ by definition, and the measured $t^{*}/z=1.00$
across all configurations confirms the implementation behaves as
specified. Tightening costs real pool: the mean stopping overcount is
$3.7z$ at $\varepsilon=0.5$ (worst-case factor $2$), $6.0z$ at
$0.25$, and $6.9z$ at $0.1$. Halving the worst-case guarantee from $3$
to $2$ therefore costs about $3.7\times$ the pool and proportionally
more Phase-2 time; the default $O=z$ remains the right operating
point, with the adaptive rule available when a tighter a priori bound
is worth paying for.

\begin{table}[t]
\caption{Pool-size settings across 72 Diabetes and Covertype
configurations. Gap is achieved-quality gap relative to the best
quality over all tested $O$; the bound column is the proven worst-case
containment ratio for the pool-optimal subset.}
\label{tab:osettings}
\centering
\begin{tabular}{lcccc}
\toprule
Setting & Within 5\% & Mean gap & Max gap & Bound \\
\midrule
$O=z$ & 53/72 & $+3.7\%$ & $+11.8\%$ & $3\cdot\OPT$ \\
$O=2z$ & 69/72 & $+1.3\%$ & $+8.2\%$ & $(1+2{\cdot}0.76)\cdot\OPT$ \\
Adaptive ($\varepsilon{=}0.5$, $O{\approx}3.7z$) & \textbf{72/72}
  & ${<}1\%$ & ${<}3\%$ & $2\cdot\OPT$ \\
\bottomrule
\end{tabular}
\end{table}

Table~\ref{tab:osettings} compares the settings. $O=z$ lands within
$5\%$ of the best achievable quality in 53/72 configurations; $O=2z$
raises that to 69/72 at exactly $2.0\times$ the runtime, buying a
further $1$--$3\%$ of quality; the adaptive rule at $\varepsilon=0.5$
reaches 72/72 with the strongest worst-case bound. We use $O=z$ as the
default and regard $O=2z$ as the conservative choice when quality
matters more than a factor of two in time.

\paragraph{What predicts the knee.}
Pearson correlation of $O_{\mathrm{knee}}/z$ against instance
statistics ($n=72$) identifies $z/k$ as the strongest structural
predictor ($r=-0.60$, $p<10^{-9}$), with raw $z$ ($r=-0.61$) and $N$
($r=-0.59$) close behind and the sparsity ratio weaker ($r=-0.43$,
$p<10^{-3}$). Dimensionality is uncorrelated ($r=+0.04$, $p=0.72$).
The three leading predictors are heavily collinear on this grid; a
joint OLS fit ($R^2=0.43$) retains $z/k$ as the strongest
standardised coefficient ($-0.25$, against $-0.13$ for raw $z$ and
$-0.19$ for $N$), while the sparsity ratio's coefficient flips sign
under the joint fit, marking its marginal correlation as a collinear
shadow; dimensionality stays null in both analyses. When each center
is responsible for many outliers, the knee hugs $z$; when outliers are
few per center, slightly more pool helps. Figure~\ref{fig:ovszk} plots
the relationship.

\begin{figure}[t]
\centering
\includegraphics[width=0.95\linewidth]{fig_synthetic.png}
\caption{Controlled synthetic sweeps. Varying $z/k$ from 2 to 120
moves $O_{\mathrm{knee}}/z$ within $[1.0,2.5]$; varying separation
from 0.5 to 8 moves it between 1.5 (tight clusters) and 1.0
(well-separated); dimensionality below 8 leaves the knee at
$O_{\mathrm{knee}}/z\approx 1.0$ with a mild rise to ${\approx}2.5$
only above $d=12$. The knee stays anchored near $z$ throughout.}
\label{fig:synthetic}
\end{figure}

Synthetic data with controlled structure confirms the picture
(Figure~\ref{fig:synthetic}): across two decades of $z/k$, a sweep of
cluster separations, and dimensions 2 through 20, the knee never
strays far from $z$, and only high dimension ($d\ge 12$) pushes it
toward $2.5z$.

\subsection{Cross-dataset validation}
\label{sec:multidata}

\begin{figure}[t]
\centering
\includegraphics[width=0.95\linewidth]{fig_ratio_both.png}
\caption{LP ratio versus $N$ on Diabetes (left) and Covertype
(right). The ordering and the magnitudes replicate the Adult results
of Figure~\ref{fig:ratiovsN}: OC-Local tracks the LP bound most
closely at every size on both datasets.}
\label{fig:ratioboth}
\end{figure}

\begin{table}[t]
\caption{LP ratios on the secondary datasets (36 and 30
configurations), against an exact LP. Charikar is reported at its
achieved solution radius.}
\label{tab:multidata}
\centering
\begin{tabular}{llcccc}
\toprule
Dataset & Algorithm & Mean & Std & Min & Max \\
\midrule
Diabetes ($d{=}8$, $n{=}36$) & RKC-Gonzalez & 1.923 & 0.256 & 1.449 & 2.414 \\
 & Ding & 1.227 & 0.048 & 1.143 & 1.330 \\
 & Charikar (achieved) & 1.201 & 0.044 & 1.145 & 1.313 \\
 & OC-Local & \textbf{1.155} & 0.043 & 1.093 & 1.275 \\
Covertype ($d{=}10$, $n{=}30$) & RKC-Gonzalez & 1.671 & 0.173 & 1.418 & 1.960 \\
 & Ding & 1.213 & 0.049 & 1.121 & 1.300 \\
 & Charikar (achieved) & 1.227 & 0.060 & 1.134 & 1.374 \\
 & OC-Local & \textbf{1.171} & 0.051 & 1.095 & 1.294 \\
\bottomrule
\end{tabular}
\end{table}

Table~\ref{tab:multidata} reports the secondary datasets. OC-Local has
the best mean LP ratio of every algorithm with a logged solution
radius: $1.155$ on Diabetes and $1.171$ on Covertype, against
Charikar's achieved $1.201$ and $1.227$, Ding's $1.227$ and $1.213$,
and Gonzalez's $1.923$ and $1.671$. The ordering of the algorithms
replicates on every dataset. The plain-Gonzalez baseline degrades
sharply on these datasets ($1.92$ and $1.67$), which widens rather
than narrows OC-Local's margin. Figure~\ref{fig:ratioboth} shows the
per-$N$ behaviour replicating the Adult curves. On HTRU2 the raw radii
tell the same story at $N=17{,}000$, $k=10$, $z=10\%$: OC-Local $1.79$
in $17$ seconds against Gonzalez $4.21$, with Charikar's achieved
radius above OC-Local's at $56$ seconds.

\begin{table}[t]
\caption{OC-Local speedup over Ding at its 99\%-certificate trial
count. OC-Local is faster in all 66 configurations.}
\label{tab:speedup}
\centering
\begin{tabular}{lcccc}
\toprule
Dataset & Mean & Median & Min & Max \\
\midrule
Diabetes ($n{=}36$) & $2.0\times$ & $2.0\times$ & $1.3\times$ & $2.4\times$ \\
Covertype ($n{=}30$) & $2.6\times$ & $2.6\times$ & $2.1\times$ & $3.6\times$ \\
\bottomrule
\end{tabular}
\end{table}

The speed advantage over Ding on these datasets is consistent but
modest (Table~\ref{tab:speedup}): $2$--$3\times$ at Ding's full
certificate trial count, in 66 of 66 configurations. At the hardest
configuration ($N=20{,}000$, $k=50$, $z=20\%$ on Diabetes) Ding takes
$639$ seconds against OC-Local's $295$. The matched Adult
configuration costs OC-Local only $57$ seconds; the median advantage
over Ding is therefore $8\times$ on Adult and $2$--$3\times$ here. The
structural difference remains---Ding's trial count grows linearly in
$z$ while OC-Local makes one pass---but each Ding trial is cheap, so
the measured gap is a constant factor, not orders of magnitude. The
decisive differences at scale are quality (Table~\ref{tab:multidata}),
determinism, and the single-pass anytime certificate; speed is a
consistent but secondary advantage.

\subsection{Scalability}
\label{sec:scalability}

\begin{table}[t]
\caption{OC-Local at large $N$ (Diabetes, $k=20$, $z=10\%$), measured.}
\label{tab:scale}
\centering
\begin{tabular}{rccc}
\toprule
$N$ & OC-Local radius & OC-Local time & Ding time \\
\midrule
20{,}000 & 2.060 & 57s & 133s \\
50{,}000 & 2.110 & 310s & 900s \\
99{,}000 & 2.150 & 1{,}100s & 3{,}800s \\
\bottomrule
\end{tabular}
\end{table}

Table~\ref{tab:scale} pushes to the full Diabetes dataset. At
$N=99{,}000$ OC-Local finishes in $1{,}100$ seconds, deterministic and
single-pass, against $3{,}800$ seconds for Ding's randomised
certificate configuration; Charikar's $O(N^{2.12})$ scaling implies
several hours, and the LP is out of reach entirely. At $10^{5}$ points
OC-Local is the fastest method that carries any formal statement about
its search space, and the only deterministic one.

\section{Conclusion}

OverclusterDRG fills the gap between quality-focused and
speed-focused algorithms for Robust $k$-Center with Outliers: a
deterministic, LP-free, single-pass method whose pool provably
contains a $3\cdot\OPT$ solution exactly when the overcount reaches
the outlier budget, whose practical variant attains a mean LP ratio of
$1.147$ under an exact LP on Adult and $1.15$--$1.17$ on two further
datasets while beating Charikar in $33/36$ and Ding in $35/36$
configurations at $7$--$8\times$ median speedups, and which remains the
best-quality, fastest deterministic method at $N\approx 10^{5}$, where
the LP is unavailable. Four problems remain open: identifying the
structural hypothesis under which the $(1+2/k)$ locality gap holds and
proving it, which demands a new technique for max-type objectives in
place of the $k$-median telescoping argument; closing the
containment-constant interval $[2,3]$ (Theorems~\ref{thm:sufficiency}
and~\ref{thm:lowerbound2})---our constructions past $2$ are repeatedly
defeated by farthest-first self-correction, suggesting the truth sits
below $3$; deriving a closed-form optimal overcount
$O^{*}=g(k,z,d,N)$ from the empirical knee at ${\approx}1.5z$ with
$z/k$ as its strongest predictor; and extending the phase-transition
analysis to $k$-median with outliers, where the sum-type objective
changes the geometry of the argument entirely.

\appendix
\section{Phase-2 Subroutines}
\label{app:pseudocode}

\begin{algorithm}[t]
\caption{OC-Forward$(X, Q, k, z)$}
\label{alg:ocforward}
\begin{algorithmic}[1]
\State Set $C = \{q_1\}$, the first pool point.
\State For each remaining pool point $q \in Q \setminus C$, compute
       $r_z(C \cup \{q\})$ --- the robust radius if $q$ is added.
\State Add to $C$ the point $q^* = \arg\min_{q \in Q \setminus C}\,
       r_z(C \cup \{q\})$.
\State Repeat Steps 2--3 until $|C| = k$.
\State \Return $C$
\end{algorithmic}
\end{algorithm}

\begin{algorithm}[t]
\caption{OC-Local$(X, Q, C, z)$}
\label{alg:oclocal}
\begin{algorithmic}[1]
\State For each center $c_i \in C$ and each unused pool point
       $q \in Q \setminus C$, compute $r_z(C - \{c_i\} + \{q\})$.
\State If any such swap strictly reduces $r_z(C)$, apply the first
       such swap found: set $C \leftarrow C - \{c_i\} + \{q\}$ and
       restart Step 1.
\State If no improving swap exists, stop.
\State \Return $C$
\end{algorithmic}
\end{algorithm}

\bibliographystyle{siamplain}
\bibliography{refs}
\end{document}
